% overload.sty の自己検査つき試験。
% 期待と異なる結果が1つでもあれば、最後に LaTeX エラーを出して失敗させる。
% 実行例: (リポジトリ直下で) TEXINPUTS=.: pdflatex -interaction=nonstopmode -halt-on-error test/overload-test.tex
\documentclass{article}
\usepackage{overload}
\makeatletter

\newcount\failcount
% #1: 試験名、#2: 実際の値を持つマクロ、#3: 期待する本体
\newcommand\expect[3]{%
  \def\@expected{#3}%
  \ifx#2\@expected
    \typeout{PASS: #1}%
  \else
    \global\advance\failcount\@ne
    \typeout{FAIL: #1: got [\meaning#2] want [\detokenize{#3}]}%
  \fi
}
% #1: 試験名、#2: 実際の箱、#3: 期待する箱(幅で比較する)
\newcommand\expectwd[3]{%
  \ifdim\wd#2=\wd#3
    \typeout{PASS: #1}%
  \else
    \global\advance\failcount\@ne
    \typeout{FAIL: #1: got \the\wd#2\space want \the\wd#3}%
  \fi
}
% #1: 試験名、#2: パッケージまたは LaTeX のエラーを起こすはずのコード
\newcommand\expecterror[2]{%
  \begingroup
    \global\let\lasterror\relax
    \def\PackageError##1##2##3{\gdef\lasterror{##2}}%
    \def\@latex@error##1##2{\gdef\lasterror{##1}}%
    \setbox0\hbox{#2}%
    \ifx\lasterror\relax
      \global\advance\failcount\@ne
      \typeout{FAIL: #1: no error raised}%
    \else
      \ifdim\wd0=\z@
        \typeout{PASS: #1}%
      \else
        \global\advance\failcount\@ne
        \typeout{FAIL: #1: error raised but stray output of width \the\wd0}%
      \fi
    \fi
  \endgroup
}

% --- 基本の分岐。名前は {\hoge} と {hoge} のどちらでも書ける ---
\newoverload{\hoge}
\addoverload{\hoge}{0}{\gdef\result{0}}
\addoverload{hoge}{1}{\gdef\result{1:#1}}
\addoverload{\hoge}{2}{\gdef\result{2:#1|#2}}

% --- 本体が引数を2回使う場合と、引数を実行する場合 ---
\newoverload{pair}
\addoverload{pair}{1}{[1:#1]}
\addoverload{pair}{2}{[2:#1|#2|#1]}
\newoverload{\raw}
\addoverload{\raw}{2}{\xdef\result{\detokenize{#1|#2|#1}}}
\newoverload{\exec}
\addoverload{\exec}{1}{#1}

% --- 登録に抜けがある場合 ---
\newoverload{\gap}
\addoverload{\gap}{1}{\gdef\result{g1}}
\addoverload{\gap}{3}{\gdef\result{g3}}

% --- 本体が別のオーバーロード呼び出しで終わる場合と、引数9個 ---
\newoverload{\fwd}
\addoverload{\fwd}{1}{\pair{#1}}
\newoverload{\nine}
\addoverload{\nine}{9}{\gdef\result{#1#2#3#4#5#6#7#8#9}}

% --- カーネルの内部名 \e@alloc と衝突しないこと ---
\newoverload{alloc}
\addoverload{alloc}{0}{}
\newcount\aftercount

\begin{document}

\hoge \expect{0 args}\result{0}
\hoge{A}\expect{1 arg}\result{1:A}
\hoge{A}{B}\expect{2 args}\result{2:A|B}
\hoge{}\expect{empty arg}\result{1:}
\hoge{a\par b}\expect{par in arg}\result{1:a\par b}

% 最大個数を超えるグループは取り込まない
\setbox0\hbox{\hoge{A}{B}{C}}\setbox2\hbox{C}
\expect{stop at max: dispatch}\result{2:A|B}
\expectwd{stop at max: rest kept}{0}{2}

% 空白で収集が止まり、空白は残る
\setbox0\hbox{x\hoge{A} {B}}\setbox2\hbox{x B}
\expect{space stops: dispatch}\result{1:A}
\expectwd{space stops: space kept}{0}{2}

% \enddelim による明示的な終端と、旧来の書き方
\setbox0\hbox{\hoge{A}\enddelim{B}}\setbox2\hbox{B}
\expect{enddelim: dispatch}\result{1:A}
\expectwd{enddelim: rest kept}{0}{2}
\hoge\enddelim \expect{legacy 0}\result{0}
\hoge{A}{B}\enddelim \expect{legacy 2}\result{2:A|B}

% 入れ子の呼び出しで引数が壊れない
\setbox0\hbox{\pair{\pair{A}}{B}}\setbox2\hbox{[2:[1:A]|B|[1:A]]}
\expectwd{nested call}{0}{2}
\raw{\pair{A}}{B}
\edef\want{\detokenize{\pair{A}|B|\pair{A}}}
\ifx\result\want \typeout{PASS: args passed as tokens}\else
  \global\advance\failcount\@ne
  \typeout{FAIL: args passed as tokens: got [\meaning\result]}\fi

% 本体末尾の呼び出しが、元の呼び出しの後ろのグループを奪わない
\setbox0\hbox{\fwd{A}{B}}\setbox2\hbox{[1:A]B}
\expectwd{forwarding body stops}{0}{2}

% 引数9個。10個目は取り込まない
\setbox0\hbox{\nine{a}{b}{c}{d}{e}{f}{g}{h}{i}{j}}\setbox2\hbox{j}
\expect{9 args}\result{abcdefghi}
\expectwd{9 args: 10th kept}{0}{2}

% 引数に # を含められる
\exec{\def\x#1{<#1>}\xdef\result{\x{q}}}
\expect{hash in arg}\result{<q>}

% \edef では展開されずに残る
\edef\moved{\hoge{A}}
\edef\want{\noexpand\hoge{A}}
\ifx\moved\want \typeout{PASS: protected in edef}\else
  \global\advance\failcount\@ne
  \typeout{FAIL: protected in edef: got [\meaning\moved]}\fi

% 定義はグループに局所的
\begingroup
  \newoverload{\localcmd}\addoverload{\localcmd}{0}{}%
\endgroup
\@ifundefined{localcmd}{\typeout{PASS: local to group}}{%
  \global\advance\failcount\@ne \typeout{FAIL: local to group}}

% エラーになるべき呼び出しと定義
\expecterror{unregistered count (2 of 1,3)}{\gap{A}{B}}
\expecterror{unregistered count (0 of 1,3)}{\gap}
\gap{A}{B}{C}\expect{gap 3}\result{g3}
\expecterror{newoverload on defined}{\newoverload{\hoge}}
\expecterror{newoverload on relax}{\newoverload{\relax}}
\expecterror{newoverload on end-name}{\newoverload{\endfoo}}
\expecterror{renewoverload on undefined}{\renewoverload{\nosuchcommand}}
\expecterror{addoverload undeclared}{\addoverload{\undeclared}{1}{stray #1}}
\expecterror{addoverload arity 10}{\addoverload{\hoge}{10}{stray}}
\expecterror{addoverload negative}{\addoverload{\hoge}{-1}{stray}}
\expecterror{addoverload duplicate}{\addoverload{\hoge}{1}{stray #1}}

% \renewoverload は登録済みの本体を消す
\renewoverload{\gap}
\addoverload{\gap}{0}{\gdef\result{new0}}
\gap \expect{renew: new body}\result{new0}
\setbox0\hbox{\gap{A}}\setbox2\hbox{A}
\expectwd{renew: old arity gone}{0}{2}

% 見出し(目次への書き出し)と表のセルで使える
\tableofcontents
\section{Title \pair{S}{T}}
\begin{tabular}{ll}
  \pair{A} & \pair{A}{B} \\
  \pair{C}{D} & x
\end{tabular}

\ifnum\failcount>\z@
  \PackageError{overload-test}{\the\failcount\space test(s) failed}{See the FAIL lines above.}
\else
  \typeout{ALL TESTS PASSED}
\fi
\end{document}
